\documentclass[12pt]{letter}
\usepackage[a4paper, top=2.5cm, bottom=2.5cm, left=2.5cm, right=2.5cm]{geometry}
\usepackage[T1]{fontenc}
\usepackage[utf8]{inputenc}
\usepackage{url}
\usepackage{parskip}
\usepackage{lmodern}
\usepackage{cmap}

\date{8 October 2026}
\address{}
\signature{}

\begin{document}

\begin{letter}{The Editor-in-Chief \\
\textit{Applied Sciences} (ISSN 2076-3417) \\
Section: Computing and Artificial Intelligence \\
MDPI AG, Basel, Switzerland}

\opening{Dear Editor,}

We submit for your consideration the manuscript entitled ``Predicting Student Test Scores with
Reliable Score Ranges: A Leakage-Free Streaming-Rating Study of 2.2 Million Records from a
Private Kazakhstani UNT-Preparation Platform'' by Abylay Salimzhanov, Artem Bykov,
Aadm Ziro, Ruslan Bulgakov, and Aasso Ziro. The paper introduces GRACE, a streaming rating architecture that produces a calibrated
score range alongside each predicted score, updating in constant time as new student attempts
arrive.

Two practical obstacles prevent machine-learning score-prediction systems from being trusted in
operational educational platforms. Evaluation accuracy is routinely inflated by data leakage, and
point predictions give educators no indication of when a forecast should be treated with caution.
Both problems appear consistently in the published literature, and both are addressed here on a
real-world dataset of 2,208,313 practice-test records from 55,482 students collected between 2018
and 2025 by a Kazakhstani university-entrance-preparation company.

A factorial audit quantifies the joint evaluation optimism produced by two widespread
feature-engineering errors combined with random train/test splitting. On this dataset---the
largest longitudinal test-preparation corpus reported from Central Asia---the audit yields a
combined $R^2$ inflation of $+0.349$, suggesting that a substantial share of accuracy
improvements claimed in prior published work reflects evaluation artifacts rather than genuine
predictive gains. The released audit pipeline provides a reusable protocol for other educational
datasets.

GRACE (Glicko-informed Rating-Augmented Confidence Estimation) is a hierarchical streaming
rating system adapted from binary-outcome psychometrics to continuous-score regression. Its
central novelty is the dual role of the Glicko rating deviation (RD): RD simultaneously serves as
a real-time uncertainty state within the rating layer and as the pre-specified grouping variable
for Mondrian conformal quantile regression, so that prediction-interval width adapts automatically
to measured student uncertainty---widening for cold-start and long-inactive students---without any
post-hoc binning decisions. The rating layer updates in $O(1)$ time per attempt with no
retraining.

GRACE is evaluated under rolling-origin temporal evaluation across five expanding annual windows
($R^2 = 0.522 \pm 0.052$), with ablation studies, paired-bootstrap significance tests (Mean
Absolute Error improvement $-1.79$ percentage points, 95\% CI $[-1.81, -1.77]$), and per-group
coverage analysis. Score ranges achieve 0.886 empirical coverage (90\% target) across three RD strata (0.880--0.891), with mean width 48.0~pp.
In 2022, the leakage-free baseline collapses to $R^2 = 0.048$ while GRACE holds at $0.437$,
illustrating the practical value of the rating features under distributional shift.

The manuscript has not been published previously, is not under consideration for publication
elsewhere, and will not be submitted elsewhere before a decision is received from
\textit{Applied Sciences}. All authors have read and approved the submitted version. The study
uses anonymized administrative records provided under a data-use agreement; no human subjects
were recruited and no personally identifiable information was processed, so institutional ethical
approval was not required. The analysis code is publicly available at
\url{https://github.com/Salimzhanov/GRACE-student-prediction}.

We believe this paper fits the Computing and Artificial Intelligence section of
\textit{Applied Sciences} and will be of interest to readers working on educational data mining,
adaptive learning systems, applied machine learning, and evaluation methodology. We would be
grateful for the opportunity to have this work reviewed.

\closing{Yours sincerely,}

\noindent Abylay Salimzhanov (Corresponding Author)\\
School of Information Technologies\\
Kazakh-British Technical University (KBTU)\\
59 Tole bi Street, Almaty 050000, Kazakhstan\\
\texttt{a.salimzhanov@kbtu.kz}\\
ORCID: 0000-0001-6630-585X

\bigskip

\noindent Artem Bykov\\
International Information Technology University (IITU)\\
34A Manas Street, Almaty 050000, Kazakhstan\\
\texttt{a.bykov@iitu.edu.kz}\\
ORCID: 0000-0002-9563-5185

\bigskip

\noindent Aadm Ziro\\
International Information Technology University (IITU)\\
34A Manas Street, Almaty 050000, Kazakhstan\\
\texttt{aziro@iitu.edu.kz}\\
ORCID: 0009-0006-7445-201X

\bigskip

\noindent Ruslan Bulgakov\\
Al-Farabi Kazakh National University (KazNU), Almaty, Kazakhstan\\
\texttt{bulgakov.ruslan.kaznu@gmail.com}\\
ORCID: 0009-0002-2036-6651

\bigskip

\noindent Aasso Ziro\\
Kazakh-British Technical University (KBTU)\\
59 Tole bi Street, Almaty 050000, Kazakhstan\\
\texttt{aa.ziro@kbtu.kz}\\
ORCID: 0000-0002-5952-877X

\end{letter}

\end{document}
